\nomenclature[00]{$\mathcal{X}_{(m)}$}{пространство наблюдений модальности~$m$}
\nomenclature[01]{$\mathcal{X} = \mathcal{X}_{(1)} \times \cdots \times \mathcal{X}_{(M)}$}{совместное пространство нейробиологических данных}
\nomenclature[02]{$\mathcal{X}_f, \mathcal{X}_e, \mathcal{X}_s$}{пространства данных фМРТ, ЭЭГ и~временных рядов фМРТ}
\nomenclature[03]{$\mathcal{Y}$}{пространство стимулов}
\nomenclature[04]{$\mathbf{F}_i, \mathbf{E}_i, \mathbf{I}_i$}{фМРТ, ЭЭГ и~изображение $i$-го наблюдения}
\nomenclature[05]{$\mathbf{S}, \mathcal{T}, L$}{временной ряд фМРТ, его длина и~число категорий стимулов}
\nomenclature[06]{$\mathbf{u}$}{бинарный ряд стимуляции}
\nomenclature[07]{$\mathfrak{D}$}{обучающая выборка}
\nomenclature[08]{$N, N_y$}{объем выборки и~число объектов класса~$y$}
\nomenclature[09]{$\mathbf{g}$}{общее отображение декодирования $\mathcal{X} \to \mathcal{Y}$}
\nomenclature[10]{$\mathbf{q}$}{прогноз фМРТ-отклика по~стимулу $\mathcal{Y} \to \mathcal{X}_f$}
\nomenclature[11]{$\mathbf{h}$}{классификатор временных рядов $\mathcal{X}_s \to \{1, \ldots, L\}$}
\nomenclature[12]{$\mathbf{f}$}{целевое отображение $\mathcal{X}_f \times \mathcal{X}_e \to \mathcal{Y}$}
\nomenclature[13]{$\Delta t$}{гемодинамический лаг BOLD-сигнала}
\nomenclature[14]{$\mu, \nu$}{частоты регистрации фМРТ и~кадров видеоряда}
\nomenclature[15]{$\boldsymbol{\delta}(t)$}{приращение фМРТ-скана между соседними кадрами}
\nomenclature[16]{$\lambda$}{коэффициент $L_2$-регуляризации}
\nomenclature[17]{$\mathcal{M}, \mathcal{M}^{y}$}{бинарная маска активности (общая и~для класса~$y$)}
\nomenclature[18]{$k_s, h$}{размер ядра сжатия и~число активных вокселей маски}
\nomenclature[19]{$\mathbf{R}_i \in \mathbb{S}_{+}^{h}$}{ковариационная матрица активных вокселей $i$-го объекта}
\nomenclature[20]{$\bar{\mathbf{R}}$}{риманово среднее ковариационных матриц}
\nomenclature[21]{$\mathbf{P}_i$}{проекция ковариационной матрицы на~касательное пространство}
\nomenclature[22]{$\mathcal{E}_f, \mathcal{E}_e, \mathcal{F}$}{энкодеры фМРТ, ЭЭГ и~модуль объединения}
\nomenclature[23]{$\mathbf{Z}_f, \mathbf{Z}_e, \mathbf{Z}_c, \mathbf{Z}_I$}{эмбеддинги фМРТ, ЭЭГ, комбинированный и~CLIP}
\nomenclature[24]{$d_f, d_e, d$}{размерности латентных пространств}
\nomenclature[25]{$\tau$}{обучаемая температура контрастивной функции потерь}
\nomenclature[26]{$\mathcal{L}_{\text{contr}}$}{контрастивная функция потерь}
\nomenclature[27]{$\mathcal{L}_{\text{prior}}$}{функция потерь априорной диффузионной модели}
\nomenclature[28]{$B$}{размер батча}
\nomenclature[29]{$\delta_{\mathcal{R}}$}{геодезическое расстояние на~многообразии $\mathbb{S}_{+}^{h}$}
\nomenclature[30]{$\mathbf{S}_i, \mathbf{p}_i$}{отбеленный касательный вектор и~вектор признаков $i$-го объекта}
\nomenclature[31]{$\omega(\mathbf{S})$}{спектральный разброс симметричной матрицы}
\nomenclature[32]{$g(\omega), c(\rho)$}{константы искажения расстояний касательной проекцией}
\nomenclature[33]{$\mathbf{M}, \bar{\mathbf{M}}$}{маска наблюдаемости записи ЭЭГ и~ее дополнение}
\nomenclature[34]{$\mathbf{E}_{\mathrm{obs}}, \mathbf{E}_{\mathrm{miss}}$}{наблюдаемая и~пропущенная части записи ЭЭГ}
\nomenclature[35]{$\mathbf{D}$}{матрица расстояний между электродами}
\nomenclature[36]{$v_{\theta}$}{векторное поле модели восстановления пропусков}
\nomenclature[37]{$\mathbb{I}$}{единичная матрица}
